% !TeX spellcheck = en_GB
% !TeX encoding = UTF-8
% !TeX root = main.tex
% !TeX TS-program = pdflatex
% !BIB TS-program = biber
\documentclass[sigconf,natbib=false,language=english]{acmart}

\usepackage[datamodel=acmdatamodel,style=acmnumeric,backend=biber]{biblatex}
\usepackage{csquotes}
\addbibresource{references.bib}
\geometry{a4paper}

\setcopyright{rightsretained}
\copyrightyear{2026}
\acmDOI{}
\acmISBN{}
\acmConference[Information Society 2026]{Information Society 2026: 29th International Multiconference}{5--9 October 2026}{Ljubljana, Slovenia}
\settopmatter{printccs=false,printacmref=false}

% Replace bracketed items as details and results are confirmed.
\newcommand{\placeholder}[1]{\textcolor{red}{[\textsc{To complete:} #1]}}
\newcommand{\citationneeded}{\textcolor{red}{[\textsc{Citation needed}]}}

\begin{document}

\title[Comparing Eye-Tracking Representations]{Raw Signals, Handcrafted Features, or Learned Embeddings? A Controlled Comparison of Eye-Tracking Representations for User-State Prediction}

% Add each author only after names, affiliations, and email addresses are final.
\author{\placeholder{Author name}}
\affiliation{%
  \institution{\placeholder{Institution}}
  \city{\placeholder{City}}
  \country{\placeholder{Country}}}
% Replace this syntactically valid address with the corresponding author's email.
\email{author@example.org}
\renewcommand{\shortauthors}{\placeholder{Author surname} et al.}

\begin{abstract}
Eye tracking can provide behavioural evidence about user state, but it remains
unclear which representation is most useful when predictions must generalise to
unseen people. This paper compares four representations of Tobii eye-tracking
windows from the TrustMe study: raw gaze and pupil signals, handcrafted
features, GazeMAE embeddings, and MOMENT embeddings. We frame the task as
predicting the provisional self-report target \textit{q5} and evaluate a
majority baseline, Random Forest, and multilayer perceptron under a
leave-one-subject-out protocol with fold-local preprocessing. We report
accuracy, balanced accuracy, macro-F1, and AUC, retaining undefined AUC values
for one-class held-out folds. \placeholder{Insert the common cohort size, the
principal quantitative result, and the resulting practical recommendation.}
\end{abstract}

\keywords{eye tracking, user-state prediction, representation learning, GazeMAE, MOMENT, subject generalisation}

\maketitle

\section{Introduction}

Eye tracking offers a detailed record of gaze behaviour and pupil dynamics that
can support inferences about user state in interactive systems
\citationneeded. Yet, choosing how to represent those signals remains a
practical and methodological decision. Raw windows retain temporal detail;
handcrafted features encode established domain knowledge and can be easier to
interpret; and learned embeddings may transfer useful structure from gaze-
specific or general time-series pretraining.

This study makes that choice explicit through a controlled comparison on Tobii
windows from the TrustMe study. The same labelled windows and subject-
generalisation protocol are used to compare raw signals, handcrafted features,
GazeMAE embeddings, and MOMENT embeddings. The provisional primary target is
\textit{q5}; its wording, scale, and substantive interpretation are documented
in Section~\ref{sec:data} before the final submission.

Our research question is: \emph{Which eye-tracking representation best
predicts the selected TrustMe self-report state for unseen subjects?} We make
three contributions: (1) a like-for-like comparison of four representation
families; (2) a leakage-safe leave-one-subject-out evaluation with fold-local
preprocessing; and (3) evidence-based guidance on the performance,
interpretability, and implementation trade-offs of each representation.

\section{Related Work}

\placeholder{Review eye tracking for user-state, engagement, cognitive-load,
and affect prediction; introduce classical gaze and pupil features; and cite
the original GazeMAE and MOMENT papers. Position this work as an empirical
comparison, not a new representation-learning method.}

Subject-dependent behavioural data can give overly optimistic results under
random window-level splits. \placeholder{Add literature supporting
subject-disjoint evaluation and fold-local normalisation.}

\section{Dataset and Preprocessing}
\label{sec:data}

\subsection{TrustMe Tobii Data}

The analysis starts from already computed representations of TrustMe Tobii
recordings. Raw participant data, derived data, model weights, and results are
not included in this repository. \placeholder{State the study purpose, ethics
or consent information where applicable, participant count, recording setup,
sampling characteristics, session structure, and exclusions.}

\subsection{Windows and Target}

Representations are aligned through a common window identifier
(\texttt{window\_uid}). Supervised analyses use labelled windows only and use
their common intersection across the four representations, ensuring that a
representation is not advantaged by a different cohort. \placeholder{Specify
window duration, overlap, ordering within participant, the exact q5 question,
response scale, binarisation rule, and final counts of subjects and windows.}

\subsection{Data Quality and Fold-Local Processing}

Raw Tobii channels comprise left and right pupil size, horizontal and vertical
gaze coordinates, and average distance. Validity fields are used during raw
signal processing. For the main protocol, all data-dependent operations,
including imputation, feature filtering, scaling where enabled, and target
centring, are fitted using training subjects only and then applied to the
held-out subject. \placeholder{Report the final missing-data policy and
quality-control exclusions.}

\section{Eye-Tracking Representations}

We compare the following families in a fixed order throughout the paper.

\subsection{Raw Signals}

Each raw window contains five channels: left and right pupil size, gaze point
coordinates $(x,y)$, and average distance. The current experimental
configuration uses 120 samples per channel with truncate-or-pad length handling.
\placeholder{Confirm the final preprocessing, input dimensionality, and
whether this configuration is retained for the submitted experiment.}

\subsection{Handcrafted Features}

The handcrafted representation is loaded from the TrustMe feature export.
Structural metadata are excluded; features with excessive missingness and
highly correlated features are filtered within each training fold.
\placeholder{List the retained feature families and the final number of
features.}

\subsection{GazeMAE Embeddings}

GazeMAE provides a gaze-specific learned embedding for each window
\cite{bautista2020gazemae}.
\placeholder{Document the checkpoint, its input signals, pooling or embedding
selection, quality-control requirements, and final dimensionality. Cite the
source paper.}

\subsection{MOMENT Embeddings}

MOMENT provides a general-purpose time-series embedding of each window
\cite{goswami2024moment}.
\placeholder{Document the pretrained model, input channels, pooling method,
quality-control requirements, final dimensionality, and any window-count
difference relative to GazeMAE. Cite the source paper.}

\section{Experimental Setup}

\subsection{Models}

We use a most-frequent-class classifier as the majority reference baseline,
a class-balanced Random Forest with 300 trees, and a multilayer perceptron.
The current MLP configuration has hidden layers of sizes 1028, 512, 256, 128,
and 64, GELU activations, Adam optimisation, early stopping, and a fixed seed.
\placeholder{Confirm hyperparameters after the final run and explain any
model-specific preprocessing.}

\subsection{Primary Evaluation Protocol}

The primary analysis is leave-one-subject-out (LOSO) evaluation with the
normalised label rule configured in the project. Each fold holds out one
participant, preventing information from that person's windows from entering
training or preprocessing. The current q5 rule centres targets using training
subjects and handles an unseen subject using the global training mean, with a
threshold of 3.0. \placeholder{Verify the interpretation of this rule and
state the exact final target definition in reader-facing language.}

\subsection{Metrics and Reporting}

We report accuracy and balanced accuracy with their deltas over the majority
baseline, macro-F1, and AUC. One-class held-out folds are retained; metrics that
are undefined for such a fold are recorded as \texttt{NaN} rather than silently
discarded. Primary summaries aggregate at the subject level.
\placeholder{Specify uncertainty intervals or paired statistical comparisons,
if used.}

\section{Results}
\label{sec:results}

\placeholder{Do not draft conclusions here before the final experiment run.
Insert a main table with representation, model, accuracy, $\Delta$ accuracy,
balanced accuracy, $\Delta$ balanced accuracy, macro-F1, AUC, and
subject-level uncertainty. Order representations as raw, handcrafted,
GazeMAE, and MOMENT; show the majority baseline first.}

\subsection{Main Representation Comparison}

\placeholder{State the principal LOSO result, compare it with the majority
baseline, and distinguish representation effects from classifier effects.}

\subsection{Subject-Level Results}

\placeholder{Add a per-subject performance figure or distribution and
row-normalised confusion matrices for the selected comparisons. Use a fixed
colour scale from 0 to 1 and discuss one-class folds.}

\section{Discussion}

\placeholder{Answer the research question using the final results. Discuss
whether learned embeddings improve on raw signals or handcrafted features;
separate performance from preprocessing cost, interpretability, and deployment
constraints; and avoid treating two-dimensional projections as proof of class
separability.}

\subsection{Limitations}

The conclusions are limited to one dataset, one windowing configuration, a
provisional self-report target, and the selected pretrained checkpoints.
\placeholder{Add final limitations concerning sample size, label noise,
quality-control decisions, and generalisability.}

\section{Conclusion}

\placeholder{Summarise the final controlled comparison, its strongest result,
and a concise recommendation for practitioners choosing an eye-tracking
representation.}

\begin{acks}
\placeholder{Add funding, institutional, and data-collection acknowledgments
only if required and approved.}
\end{acks}

\printbibliography

\appendix
\section{Reproducibility Details}

\placeholder{Add the exact YAML configuration, representation dimensions,
common-intersection counts, software versions, per-fold metrics, and secondary
protocol or projection results that do not fit in the main paper.}

\end{document}
